
Netdata es una plataforma de código abierto diseñada para supervisar, en tiempo real, el estado y rendimiento de sistemas informáticos, servidores, dispositivos y servicios de red. Su función principal es recopilar métricas del sistema y presentarlas mediante un panel gráfico interactivo, facilitando la identificación de fallos, sobrecargas y comportamientos anormales.

Entre las métricas que Netdata puede monitorear se encuentran el uso del procesador, la memoria RAM, el almacenamiento, la temperatura del equipo, los procesos en ejecución y el tráfico de las interfaces de red. También permite observar el funcionamiento de diferentes aplicaciones y servicios, como servidores web, bases de datos, contenedores y otros componentes de infraestructura. \cite{netdata-documentacion}

Netdata trabaja mediante un agente que se instala en el dispositivo que se desea supervisar. Este agente recopila las métricas disponibles y las muestra a través de una interfaz web. La información se actualiza con alta frecuencia, lo que permite observar prácticamente en tiempo real los cambios que ocurren en el sistema. Además, el software puede generar alertas cuando una métrica supera determinados límites o presenta un comportamiento fuera de lo normal.

Una de sus principales características es que puede detectar automáticamente una gran cantidad de componentes y servicios, reduciendo la configuración manual necesaria durante su instalación. Asimismo, puede utilizarse de manera independiente dentro de una red local o integrarse con Netdata Cloud para administrar y visualizar varios equipos desde una interfaz centralizada. \cite{netdata-documentacion}

